The TRADES arm is small: \tradesSeeds{} seeds, each with an unperturbed twin and five interventions (\nTradesEpisodes{} episodes of about \tradesWall{} minutes each). The parent order is \tradesQty{} shares, 10\% of the generated market's mean volume in the 10-minute window. Twin and intervention were identical before the first child order in \tradesPreIdentical{} of runs and differed afterwards in all of them, so the generated flow does respond to the injected orders.

\begin{table*}[!htbp]
\centering\footnotesize
\caption{The 10-minute task (sell 10\% of background volume in 20 slices) in every ABIDES world and in the retrained TRADES market T1. Paired impact cost in bps, mean $\pm$ 95\% half-width. Bold: cheapest schedule in each column.}\label{tab:short}
\begin{tabular}{lrrrrrrr}
\toprule
Policy & A1 & A2 & A3 & A4 & A5 & A6 & T1 \\
\midrule
Schedule $b=-2$ (back-loaded) & 0.7{\tiny$\,\pm$0.2} & \textbf{0.4{\tiny$\,\pm$0.3}} & 11.8{\tiny$\,\pm$9.1} & 1.0{\tiny$\,\pm$1.1} & \textbf{0.5{\tiny$\,\pm$0.2}} & 1.2{\tiny$\,\pm$1.2} & 41.6{\tiny$\,\pm$21.2} \\
Schedule $b=0$ (TWAP) & 0.7{\tiny$\,\pm$0.3} & 0.4{\tiny$\,\pm$0.3} & 4.6{\tiny$\,\pm$10.8} & 1.2{\tiny$\,\pm$1.2} & 0.6{\tiny$\,\pm$0.2} & 0.8{\tiny$\,\pm$0.7} & 50.6{\tiny$\,\pm$27.2} \\
Schedule $b=2$ & \textbf{0.5{\tiny$\,\pm$0.3}} & 0.6{\tiny$\,\pm$0.3} & 4.4{\tiny$\,\pm$6.3} & \textbf{0.7{\tiny$\,\pm$0.6}} & 0.6{\tiny$\,\pm$0.1} & \textbf{0.3{\tiny$\,\pm$0.6}} & 44.2{\tiny$\,\pm$28.8} \\
Schedule $b=8$ (front-loaded) & 1.1{\tiny$\,\pm$0.3} & 1.2{\tiny$\,\pm$0.3} & \textbf{3.8{\tiny$\,\pm$2.3}} & 1.5{\tiny$\,\pm$0.6} & 0.7{\tiny$\,\pm$0.1} & 1.6{\tiny$\,\pm$0.6} & \textbf{26.6{\tiny$\,\pm$16.6}} \\
Immediate & 2.1{\tiny$\,\pm$0.4} & 1.9{\tiny$\,\pm$0.5} & 5.3{\tiny$\,\pm$2.3} & 2.5{\tiny$\,\pm$0.9} & 1.1{\tiny$\,\pm$0.2} & 4.5{\tiny$\,\pm$1.4} & 15.7{\tiny$\,\pm$2.4} \\
\midrule
Episodes per cell & 40 & 40 & 40 & 40 & 40 & 40 & 6 \\
\bottomrule
\end{tabular}
\end{table*}


Table~\ref{tab:short} puts the TRADES market next to the ABIDES worlds on the identical task. The two families give opposite advice. In every stable ABIDES world the immediate order is clearly the most expensive policy and the schedules are close to one another. In the TRADES market the ordering is reversed: the immediate order and the most front-loaded schedule are the cheapest, at \tradesCostImm{}\,bps and \tradesCostFront{}\,bps ($b=8$), and TWAP is the most expensive, at \tradesCostTwap{}\,bps. Paired by seed, $b=8$ differs from TWAP by \tradesDiffFront{}\,bps (95\% interval \tradesDiffFrontCI{}; cheaper in \tradesDiffFrontN{} seeds), and the immediate order by \tradesDiffImm{}\,bps (\tradesDiffImmCI{}). The estimates are noisy: two seeds have TWAP costs near 100\,bps. The generated market is also far more impact-sensitive than any ABIDES world: selling 10\% of volume shifts the mid by \tradesDispMin{} to \tradesDispMax{}\,bps by the deadline depending on the schedule, against at most \abidesShortDispMax{}\,bps under TWAP in the stable ABIDES worlds, and the shift under TWAP is still \tradesDispTwapPost{}\,bps two minutes later. Selecting the schedule in the TRADES market and deploying it in a stable ABIDES world costs between \tradesToAbidesMin{} and \tradesToAbidesMax{}\,bps of regret; selecting in a stable ABIDES world and deploying in the TRADES market costs between \abidesToTradesMin{} and \abidesToTradesMax{}\,bps. These are point estimates from \tradesSeeds{} seeds with no train/test split.

So here is a pair of simulators that do disagree materially. The disagreement, however, has a mundane source, and it is one that realism statistics catch. The retrained generator does not hold the market it inherits. Across unperturbed twins the quoted spread is \tradesSpreadReplay{}\,bps in the last minute of real replay, \tradesSpreadFirst{}\,bps in the first generated minute, \tradesSpreadMid{}\,bps after five minutes and \tradesSpreadLast{}\,bps in the last minute of the task. A seller who trades early crosses a spread close to the real one; a seller who waits crosses a spread several times wider. Front-loading wins because the simulated market deteriorates with time since hand-over, not because of anything about impact. Table~\ref{tab:tradesreal} shows the same thing in aggregate: over the task window the generated market's spread and volatility are well above those of the real day it was trained on.

\begin{table}[!htbp]
\centering\small
\caption{Unperturbed retrained-TRADES market against the real day it was trained on, over the task window. The real window lies in the training split.}\label{tab:tradesreal}
\begin{tabular}{lrrr}
\toprule
Statistic (09:46--09:56) & TRADES mean & TRADES s.d. & Real INTC \\
\midrule
Quoted spread (bps) & 24.06 & 13.23 & 3.66 \\
10\,s return volatility (bps) & 7.29 & 1.25 & 3.92 \\
Return autocorrelation, lag 1 & -0.38 & 0.09 & -0.15 \\
$|$return$|$ autocorrelation, lag 1 & 0.23 & 0.08 & -0.15 \\
Traded volume (shares) & 315618 & 39337 & 247627 \\
\bottomrule
\end{tabular}
\end{table}


We stress what this does and does not show. It is evidence about a TRADES architecture trained for \tradesSteps{} steps on one day, and says nothing about the published model. But it illustrates the second way, besides instability, in which a simulator can mislead an execution study: a generator whose state drifts away from its conditioning imposes a time preference of its own on every schedule evaluated inside it. A spread-versus-time plot of unperturbed runs is enough to detect it.
